\documentclass[10pt]{article}

% graphicx package, useful for including eps, pdf, and raster graphics
\usepackage{graphicx}
\DeclareGraphicsExtensions{.png,.png,.jpg,.pdf}

% basic packages
\usepackage{fancyhdr}
\usepackage{color}
\usepackage{parskip}
\usepackage{float}
\usepackage{microtype}
\usepackage{url}
\usepackage{hyperref}
\usepackage{booktabs}
\usepackage{makecell}
\usepackage{multirow}
\usepackage{pbox}
\usepackage{enumitem}

% fonts
\usepackage{fontspec}
\setmainfont{DroidSans}[
  Extension   = .ttf,
  UprightFont = *,
  BoldFont    = *-Bold ]
\setmonofont{DroidSansMono}[
  Extension   = .ttf,
  UprightFont = * ]

\hypersetup{colorlinks,urlcolor=blue}
\urlstyle{same}

% page layout
\usepackage[top=1in, bottom=1in, left=0.75in, right=0.75in, includefoot, heightrounded]{geometry}

% Useful for seeing the layout as you tweak it.
%\usepackage{showframe}

\setlength{\headheight}{1.1in}
\pagestyle{fancy}
\fancyhf{}

\lhead{\includegraphics[height=1in,clip]{logo}}
\rhead{\pbox[b][1in][c]{\textwidth}{
    \small
    \textbf{Ang Pagsusuri ay Isinagawa Ng:}\\
    Northwest Genomics Center\\
    University of Washington School of Medicine\\
    Genome Sciences Box 355065\\
    3720 15th Ave NE\\
    Seattle, WA 98195-5065}}

\newcommand{\PageLine}{\rule{\textwidth}{0.25mm}}

\begin{document}

\begin{description}[font=\normalfont,align=left,labelwidth=12em]
\item [\textbf{\VAR{pat_name|e}}]
\item [\textbf{Petsa ng Kapanganakan}] \VAR{birth_date|e}
\item [\textbf{Specimen}] Gitnang Bahagi ng Ilong, sariling kinolekta
\item [\textbf{Barcode ng Pagkolekta}] \VAR{qrcode|e}
\item [\textbf{Petsa ng Pagkolekta}] \VAR{collect_ts|e}
\item [\textbf{Petsa ng Ulat}]
  %- if status_code == "never-tested"
  Hindi pa nasubukan
  %- else
  \VAR{result_ts|e}
  %- endif
\end{description}

\PageLine

\begin{center}
\Large
\textbf{RESULTA NG PANANALIKSIK – Hinihikayat kang ibahagi ang resultang ito sa
  iyong provider ng pangangalagang pangkalusugan}
\end{center}

\bigskip

\large \underline{\textbf{Tungkol sa Pagsusuri na Ito}}

\begin{itemize}
\item

  Ang COVID-19 ay isang \textbf{sakit} na dulot ng \textbf{virus} na tinatawag
  na SARS-CoV-2. Hinahanap ng pagsusuri ng pananaliksik na ito ang virus sa
  sample na kinolekta mula sa iyong ilong para matukoy kung mayroon kang
  COVID-19. Ang tekninal na katawagan para sa ganitong uri ng pagsusuri ay
  ``reverse-transcription polymerase chain reaction,'' o ``RT-PCR'' kung
  pinaikli.

\end{itemize}

\bigskip

\large \underline{\textbf{Mahahalagang Paalala}}

Bago tingnan ang resulta ng iyong pagsusuri ng pananaliksik, dapat mong suriin ang sumusunod:

\begin{itemize}
\item

  Ang \textbf{uri ng sample} na kinolekta mo at ang \textbf{pagsusuri sa lab} ay
  bahagi ng \textbf{pananaliksik na pag-aaral (research study)} at inilalarawan
  ng ulat na ito ang iyong \textbf{resulta ng pagsusuri ng pananaliksik.}

\item

  Kinolekta mo ang sample na ito at nilagyan mo (o para sa ibang tao) ito ng
  label at ibinalik sa lab sa pamamagitan ng serbisyo sa pag-deliver. Hindi ito
  kinolekta, nilagyan ng label, o ihinatid ng isang lisensiyadong propesyonal sa
  kalusugan tulad ng nurse o doktor. Maaaring makaapekto ito sa resulta ng iyong
  pagsusuri, at maaaring gusto kang muling ipasuri ng isang provider ng
  pangangalagang pangkalusugan.

\item

  Hindi pamalit sa iyong normal na medikal na pangangalaga ang pagiging bahagi
  ng \textbf{pananaliksik na pag-aaral} na ito. Kung sobrang masama ang iyong
  pakiramdam o kaya ay nababahala ka sa resulta ng iyong pagsusuri, dapat kang
  makipag-ugnayan sa doktor o provider ng pangangalagang pangkalusugan.

\end{itemize}

\bigskip
\PageLine

\large \underline{\textbf{Ang Resulta ng Pagsusuri ng SARS-CoV-2 (COVID-19)}}

%- if status_code == "negative"
Walang Natukoy (Negatibo)\\
%- elif status_code == "positive"
Natukoy (Positibo)\\
%- elif status_code == "inconclusive"
Hindi Tiyak\\
%- elif status_code == "never-tested"
Hindi Nasuri\\
%- else
\VAR{0/0} %# poor man's raise
%- endif

\PageLine
\bigskip

\large \underline{\textbf{Pag-intindi sa Iyong Resulta}}

Ipinapakita ng table na ito ang lahat ng mga resulta na maaaring ipakita ng
iyong pagsusuri at inilalarawan nito kung ano ang ibig sabihin ng bawat isa sa
mga resultang ito. Maaari mong i-print o i-email ang \textbf{resulta ng
pagsusuri ng pananaliksik} na ito para ibahagi sa iyong doktor o provider ng
pangangalagang pangkalusugan.

Anuman ang resulta ng iyong pagsusuri, dapat kang manatili sa bahay at dapat
mong iwasang ang ibang tao kung masama ang iyong pakiramdam. \textbf{Kung
nagpapatuloy ang mga sintomas mo, sobrang masama ang iyong pakiramdam o kaya ay
nababahala ka sa resulta ng iyong pagsusuri, o mayroon kang mga tanong tungkol
sa resulta ng pagsusuri ng pananaliksik, dapat kang makipag-ugnayan sa doktor o
provider ng pangangalagang pangkalusugan.}

\begin{itemize}


\item

  \textbf{If your test says ``Not Detected (Negative)''}, it means our lab did not find the
  virus that causes COVID-19 in the sample that you sent to us. No test is
  perfect, and there is a small chance that you could have the virus even though
  your test is negative.

\item

  \textbf{If your test says ``Detected (Positive)''}, it means our
  lab found the virus that causes COVID-19 in the sample that you sent us. You
  should stay home and avoid being around other people, especially if you are
  feeling sick.

  Your result will be shared with the Washington State and local health
  departments. You do \textbf{not} need to report these results yourself. The
  health department may call you to ask questions about your illness and people
  you may have been around while you were sick.

\item

  \textbf{If your test says ``Inconclusive''}, it means our lab is not certain
  of your test result. No test is perfect, and an inconclusive result happens
  when we find a small amount of SARS-CoV-2 in your sample, but not enough for
  us to say that you have COVID-19.

\item

  \textbf{If your test says ``Unable to be tested''}, it means there was either
  not enough material on your sample for us to test, or there was an error when
  we ran the test. Since we are not able to test your sample, we could not find
  whether or not you have the virus that causes COVID-19.

\end{itemize}

\bigskip

\large \underline{\textbf{Mga Mapagkukunan}}

Kung gusto mo ng higit pang impormasyon tungkol sa COVID-19, mangyaring tingnan
ang mga mapagkukunang ito mula sa mga partner sa pampublikong kalusugan:

\begin{itemize}
\item
  Public Health \textemdash Seattle \& King County (Pampublikong Kalusugan
  \textemdash Seattle at King County): \\
  \url{https://www.kingcounty.gov/depts/health/covid-19/languages/filipino.aspx}

\item

  Washington State Department of Health (Departamento ng Kalusugan ng Estado ng
  Washington): \\
  \url{http://www.doh.wa.gov/Emergencies/Coronavirus}

\item
  U.S. Centers for Disease Control and Prevention (CDC, Mga Sentro sa Pagkontrol
  at Pag-iwas sa Sakit sa Estados Unidos):\\
  \url{https://wwwn.cdc.gov/pubs/other-languages?Sort=Lang%3A%3Aasc&Language=Tagalog}

\end{itemize}

Kung gusto mong makipag-usap sa isang tao sa telepono, mangyaring tumawag sa isa
sa mga sumsuunod:

\begin{itemize}

\item

  \textbf{Mga medikal na tanong na nauugnay sa COVID-19:} Makipag-ugnayan sa King County Call Center
  (Sentro sa Pagtawag ng King County) sa pagitan ng 8AM at 7PM sa
  \textbf{206-477-3977}

\item

  \textbf{Mga hindi medikal na tanong tungkol sa COVID-19, kasama ang mga isyu
  na nauugnay sa pagsunod at negosyo:} Makipag-ugnayan sa King County COVID-19
  Business and Community Information Line (Linya ng Impormasyon sa COVID-19 sa
  Negosyo at Komunidad ng King County), Lunes – Biyernes sa pagitan ng 8:30AM at
  4:30PM sa \textbf{206-296-1608}

\item

  \textbf{Mga pangkalahatang tanong tungkol sa COVID-19 sa Estado ng
  Washington:} Makipag-ugnayan sa Washington State Novel Coronavirus Call Center
  (Sentro sa Pagtawag ukol sa Bagong Coronavirus sa Estado ng Washington) sa
  pagitan ng 6AM at 10PM sa \textbf{1-800-525-0127}

\end{itemize}

\bigskip

\large \underline{\textbf{Mga Paraan at Limitasyon}}

\textbf{Mga Paraan}\\
Ginawa ang pagsusuring ito para tukuyin ang SARs-CoV-2 gamit ang
reverse-transcription PCR. Gumagamit ang pagsusuring ito ng pagpaparami ng kopya
ng mga bahagi ng virus, ang ORf1B at S, at tinutukoy din nito ang kontrol sa
pagpaparami na human RnaseP. Ang pagsusuring ito ay binuo at tinukoy ang mga
katangian sa pagganap nito ng Northwest Genomics Center ng University of
Washington. Ang laboratoryo ay binigyan ng lisensiya ng Site para sa Medikal na
Pagsusuri ng WA State Department of Health bilang kwalipikadong magsagawa ng
lubos na kumplikadong pagsusuri sa klinikal na laboratoryo.

\textbf{Mga Limitasyon}

\begin{itemize}

\item

  Ibinibigay ang mga resultang ito bilang bahagi ng isang pananaliksik na
  pag-aaral.

\item

  Kinolekta mo ang sample na ito at nilagyan mo (o para sa ibang tao) ito ng
  label at ibinalik sa lab sa pamamagitan ng serbisyo sa pag-deliver. Hindi ito
  kinolekta, nilagyan ng label, o ihinatid ng isang propesyonal sa
  pangangalagang pangkalusugan tulad ng nurse o doktor. Maaaring makaapekto ito
  sa resulta ng iyong pagsusuri, at maaaring gusto kang muling ipasuri ng isang
  provider ng pangangalagang pangkalusugan.

\item

  Dahil batay ang ulat na ito sa data mula sa isang pananaliksik na pag-aaral,
  maaaring hindi matukoy ng iyong mga resulta ang COVID-19, ang sakit na dulot
  ng impeksiyon ng SARS-CoV-2.

\end{itemize}

\end{document}
